%----------------------------------------------------------------------------------------
%----------------------------------------------------------------------------------------
\doublespacing

\chapter{Metodología}
% Se asume que este es el Capítulo 3

El presente capítulo describe la metodología empleada para el desarrollo de la investigación titulada “Adaptación de la Arquitectura Hexagonal para mejorar la Portabilidad de Sistemas Robóticos en ROS 2”.

\section{Revisión Bibliográfica}

La revisión bibliográfica se realiza en el Capítulo II (Marco Teórico y Antecedentes) con el propósito de establecer las bases teóricas de la Arquitectura Hexagonal, el middleware ROS 2 y los entornos de simulación robótica. El resultado de esta revisión permite identificar las limitaciones de los enfoques arquitectónicos tradicionales respecto al acoplamiento tecnológico y justifica la necesidad de un nuevo enfoque basado en la separación de responsabilidades para garantizar la portabilidad entre simuladores dispares como Gazebo e Isaac Sim.

\section{Especificación de la Propuesta de Mejora de Conocimiento (Hipótesis)}

La propuesta plantea la adaptación de la Arquitectura Hexagonal (Patrón de Puertos y Adaptadores) como base estructural para el diseño de software en sistemas robóticos móviles. Su organización se define en tres capas principales (Dominio, Puertos y Adaptadores), promoviendo el desacoplamiento total de la lógica de navegación respecto a la infraestructura de simulación y el middleware.

En este marco, se formula la siguiente hipótesis:

\textbf{La adaptación de la Arquitectura Hexagonal en sistemas robóticos desarrollados con ROS 2 mejora significativamente la portabilidad del software al migrar entre entornos de simulación dispares, permitiendo la reutilización del código del dominio sin modificaciones.}


\section{Desarrollo del Artefacto}

El artefacto desarrollado en esta investigación consiste en un sistema de software modular para la navegación autónoma de un rover, diseñado bajo los principios de la Arquitectura Hexagonal. Este sistema establece una organización estricta donde:
\begin{itemize}
    \item El \textbf{Dominio} encapsula los algoritmos de navegación y control en C++ puro, sin dependencias de ROS 2.
    \item Los \textbf{Puertos} definen interfaces abstractas para los sensores (LIDAR, Odometría) y actuadores.
    \item Los \textbf{Adaptadores} implementan la comunicación con la infraestructura, traduciendo los mensajes de ROS 2 (provenientes de Gazebo o Isaac Sim) al formato interno del dominio.
\end{itemize}
Esta estructura garantiza que el núcleo del sistema sea agnóstico al simulador utilizado.

\section{Planeamiento de la Implementación de la Propuesta}

La implementación de la propuesta se planificará de acuerdo con la estructura arquitectónica definida para el artefacto. Esta planificación contempla las siguientes etapas:
\begin{enumerate}
    \item \textbf{Diseño de Interfaces (Puertos):} Definición formal de los contratos de comunicación entre el dominio y el exterior.
    \item \textbf{Implementación del Núcleo (Domain):} Desarrollo de la lógica de negocio aislada.
    \item \textbf{Desarrollo de Adaptadores de Infraestructura:} Creación de componentes específicos para interactuar con Gazebo y, posteriormente, con Isaac Sim.
    \item \textbf{Integración y Despliegue:} Configuración de los archivos de lanzamiento (launch files) y nodos de arranque (Bootstrap) para inyectar las dependencias adecuadas según el entorno de simulación seleccionado.
\end{enumerate}

Como parte del planeamiento, se establecerá el esquema de evaluación definiendo qué métricas de portabilidad e independencia (como la Métrica de Inestabilidad de Martin) serán examinadas para validar el desacoplamiento.

\section{Aplicación de la Propuesta de Mejora de Conocimiento}

La propuesta será aplicada mediante la construcción progresiva del sistema de navegación en un entorno de desarrollo controlado (Ubuntu/ROS 2 Humble). Durante esta aplicación, se verificará la capacidad del sistema para operar inicialmente en el simulador estándar (Gazebo Classic). Una vez validado el funcionamiento básico, se procederá a sustituir los adaptadores de infraestructura para conectar el mismo núcleo de navegación al simulador de alta fidelidad (NVIDIA Isaac Sim), sin alterar el código fuente del dominio. Este proceso demostrará la viabilidad técnica de la migración y la efectividad del desacoplamiento arquitectónico.

\section{Evaluación de Resultados}

La evaluación de los resultados se realizará analizando el comportamiento del artefacto en ambos escenarios de simulación. Se emplearán dos criterios fundamentales:
\begin{itemize}
    \item \textbf{Evaluación Estructural:} Medición del grado de desacoplamiento utilizando métricas de ingeniería de software (Acoplamiento Aferente/Eferente e Inestabilidad).
    \item \textbf{Evaluación Funcional (Sim-to-Sim):} Verificación de que el robot cumple la misión de navegación (ir del punto A al punto B evadiendo obstáculos) de manera equivalente tanto en Gazebo como en Isaac Sim.
\end{itemize}
Esta evaluación permitirá determinar objetivamente si la Arquitectura Hexagonal cumple con la hipótesis planteada de mejorar la portabilidad.

\section{Conclusiones}

Este paso consiste en sintetizar los hallazgos obtenidos tras la evaluación experimental, determinando si la adaptación de la Arquitectura Hexagonal efectivamente facilita la migración entre simuladores y reduce el esfuerzo de ingeniería asociado al cambio de tecnología. Las conclusiones permitirán validar la hipótesis de investigación y establecer recomendaciones arquitectónicas para futuros desarrollos de sistemas robóticos complejos en ROS 2.
